% ============================================
% Response to Reviewer 1
% ============================================
\section*{Response to Reviewer 1}

We thank the reviewer for their constructive feedback. Below we address each point in detail.

% ============================================
\subsection*{R1-1. Absolute biomass as an alternative explanation for Dominance}

\reviewercomment{I found the Author's explanation of coalescence generally convincing. However, I wonder if a simple, alternative explanation for many of the results could be that parent communities attain very different absolute densities. For example, if community A reaches 10x the overall density of community B, then even relatively rare species in community A might be initially more abundant than the dominant species in community B (because communities are combined in equal volumes, line 119). This might lead to the final community being dominated by species from A simply because species survive proportionally to their initial \emph{absolute} abundance. If I understand correctly, this scenario is different from the one tested in ED Fig.~5. High heterogeneity in absolute densities could potentially explain both the prevalence of Dominance outcomes and the difference in pairwise selection correlation between same- and cross-parent pairs. In principle, heterogeneity in absolute densities could also increase with nutrient concentration. To be clear, I have no reason to expect this pattern, but I wonder if the Authors could do more to rule it out. Since ODs were measured throughout (line 379), could those data be used to check this possibility?}

\responselabel
We thank the reviewer for raising this important alternative explanation. We tested it directly by asking whether parental-community OD$_{600}$, measured at the end of the seven-day assembly phase, predicts the parental dominance direction or the pairwise selection correlation.
If Dominance were primarily driven by unequal absolute biomass at mixing, the higher-OD parental community should win more frequently, and signed $\Delta\text{OD}_{A-B}$, the difference between the OD of parental community A and the OD of parental community B, should be positively associated with PDI. We observed the opposite pattern. In Base, the higher-OD parental community wins only $37\%$ of Dominance events, and the continuous relationship between signed $\Delta\text{OD}_{A-B}$ and PDI is weak. This opposite-to-prediction pattern becomes strongest in Nutr$+$, where Dominance is most frequent: the higher-OD parental community wins only $9/68 = 13.2\%$ of Dominance events ($p = 3.9 \times 10^{-10}$), and signed $\Delta\text{OD}_{A-B}$ is strongly negatively correlated with PDI ($\rho = -0.60$, $p = 4.0 \times 10^{-10}$; Response Fig.~R1-1A,B). Nutr$-$ is less informative for this directional test because both OD differences and Dominance frequency are smaller.

This density-based explanation also predicts that higher-OD parental communities should show stronger within-community pairwise selection correlations. To test this, we binned parental communities into low, mid, and high OD tertiles within each medium and computed the pairwise selection correlation among species from the same parental community in each OD bin (Response Fig.~R1-1C). The cross-community correlation is not re-binned by OD and is shown as a medium-level reference. Within-community pairwise selection correlation does not show a consistent increase with parental-community OD; instead, the pattern is not consistent across media. In Nutr$+$, the slight negative OD trend mirrors the negative signed $\Delta\text{OD}_{A-B}$--PDI association in the directional analysis above, rather than the positive trend expected if higher biomass drove winning. Overall, this analysis implies that differences in parental-community biomass do not explain Dominance direction or the nutrient-dependent pattern of correlated species fates.

We have added a short paragraph to the Results and a new Supplementary Figure stating: \mschange{``As an alternative hypothesis, nutrient concentration could change biomass differences among parental communities; if higher nutrient levels created larger biomass asymmetries, Dominance could reflect preferential survival of cells from the higher-biomass parental community rather than interaction-mediated co-selection. Under this hypothesis, higher-OD parental communities should more often win Dominance events and show stronger within-community pairwise selection correlations. Instead, the higher-OD parental community won only $37\%$ of Base Dominance events and $13.2\%$ of Nutr$+$ Dominance events (9/68; binomial $p = 3.9 \times 10^{-10}$; Supplementary Fig.~33), signed $\Delta\text{OD}_{A-B}$ was weakly related to PDI in Base and strongly negatively associated with PDI in Nutr$+$ ($\rho=-0.60$, $p=4.0 \times 10^{-10}$), and OD-binned within-community pairwise selection correlation did not increase consistently across media. These results indicate that parental-community biomass differences do not explain the nutrient-dependent increase in Dominance or correlated species fates.''}

\begin{figure}[H]
\centering
\includegraphics[width=\textwidth]{Fig_R1_1A_winner_loser_OD.pdf}
\caption*{\textbf{Response Fig.~R1-1A.} \textbf{Higher parental-community OD does not preferentially predict the winner in Dominance events.} Winner OD vs loser OD scatter is shown per medium; dashed line is $y=x$. Points above the line are cases where the winning parental community had higher OD than the losing parental community. The higher-OD parental community wins only $37\%$, $37\%$, and $13\%$ of Dominance events in Nutr$-$, Base, and Nutr$+$ respectively; the strongest opposite-to-prediction pattern occurs in Nutr$+$. The displayed $p$ value is from a two-sided binomial test against the null expectation that the higher-OD parental community wins $50\%$ of Dominance events.}
\end{figure}

\begin{figure}[H]
\centering
\includegraphics[width=\textwidth]{Fig_R1_1B_OD_vs_PDI.pdf}
\caption*{\textbf{Response Fig.~R1-1B.} \textbf{Continuous OD differences do not support the prediction that the higher-OD parental community wins.} Signed $\Delta\text{OD}_{A-B}$ ($\text{OD}_{A} - \text{OD}_{B}$) vs PDI is shown per medium. PDI is the Parental Dominance Index defined in the manuscript, $\text{PDI} = (2/\pi)\arctan(u/v)$. Spearman $\rho$ is the rank correlation between signed $\Delta\text{OD}_{A-B}$ and PDI; the displayed $p$ value is from a two-sided Spearman rank-correlation test of the null hypothesis of no monotonic association. The relationship is weak in Base and strongly negative in Nutr$+$ ($\rho = -0.60$, $p = 4 \times 10^{-10}$), opposite the sign predicted if higher OD caused a parental community to win. Gray points are reflections included for symmetry.}
\end{figure}

\begin{figure}[H]
\centering
\includegraphics[width=0.85\textwidth]{Fig_R1_1C_pairwise_corr_vs_OD.pdf}
\caption*{\textbf{Response Fig.~R1-1C.} \textbf{Higher parental-community OD does not produce a consistent increase in within-community pairwise selection correlation.} Rows: Nutr$-$, Base, Nutr$+$. \textbf{Left column:} parental communities are grouped into low, mid, and high within-medium OD tertiles; smaller text below the group labels shows the OD$_{600}$ range of each tertile. Bars show within-community pairwise selection correlation with bootstrap 95\% CIs; the dashed gray line and shaded band show the medium-level mean cross-community correlation and its bootstrap 95\% CI. \textbf{Right column:} per-event parental-community observations plotted by OD and within-community pairwise selection correlation; vertical dotted lines mark the low/mid/high OD boundaries used in the left column. Spearman $\rho$ and $p$ values are shown in each panel; p-values are from two-sided Spearman rank-correlation tests of the null hypothesis of no monotonic association between OD and within-community pairwise selection correlation.}
\end{figure}

% ============================================
\subsection*{R1-2. Does pH mismatch predict Dominance?}

\reviewercomment{It seems like an even more straightforward prediction would be that Dominance should become more likely when combining parent communities with qualitatively different pH (e.g.\ alkaline and acidic) compared to cases where both parents are either acidic or alkaline. Is this the case?}

\responselabel
Not significantly for Dominance frequency. We thank the reviewer for suggesting this direct test. Classifying parental-community pairs by pH type (threshold pH $= 7.0$) gives 88 acid--alk, 40 acid--acid, and 132 alk--alk events across 260 total; Nutr$-$ contains no acidic parental communities, so the per-medium analysis is restricted to Base and Nutr$+$. Comparing acid--alk to pooled same-pH pairs (acid--acid $\cup$ alk--alk), Dominance frequency is slightly elevated in Nutr$+$ ($79.5\%$, 35/44 vs $71.7\%$, 33/46; OR $= 1.53$, Fisher $p = 0.47$) and slightly reduced in Base ($61.4\%$ vs $69.4\%$; OR $= 0.70$, $p = 0.49$). These comparisons are not significant, and some categories are small (Base acid--acid $n = 6$; Nutr$+$ alk--alk $n = 12$).

\reviewercomment{To support the hypothesis that the top-down regime emerges because dominant species strongly modify the pH of the ecosystem, the Authors show that the predictive power of the parent community pH increases in the Nutr+ condition when looking at coalescence between communities where one parent is acidic and one is alkaline (lines 261, ED Fig.~8).}

Yes for winner identity in Nutr$+$. The directional test is clearer. Defining signed PDI $= +\text{PDI}$ if the acidic parental community won and $-\text{PDI}$ otherwise, in Nutr$+$ the acidic parental community wins $86.4\%$ (38/44) of acid--alk events (binomial $p = 9.4 \times 10^{-7}$; mean signed PDI $= +0.64$, one-sample $t$-test $p = 2.4 \times 10^{-8}$), while Base shows a weaker $+0.28$ signal (acid-wins $63.6\%$, $p = 0.10$; signed PDI $t$-test $p = 0.030$; Response Fig.~R1-2). Thus, pH mismatch is more informative as a predictor of the winner in Nutr$+$ than as a predictor of whether an event is classified as Dominance.

We added to Results \S2.5: \mschange{``Whereas pH mismatch alone did not significantly increase Dominance frequency within each medium, within acid--alk coalescence pairs the acidic parental community was the systematic winner in Nutr$+$ (38/44 events, binomial $p < 10^{-6}$; Supplementary Fig.~34), consistent with acidification contributing to winner identity under high-nutrient conditions.''}

\begin{figure}[H]
\centering
\includegraphics[width=\textwidth]{Fig_R1_2_acidalk_per_medium.pdf}
\caption*{\textbf{Response Fig.~R1-2.} \textbf{pH mismatch weakly predicts outcome class, but acidic parental communities preferentially win acid--alk events in Nutr$+$.} \textbf{(Left, Middle)} Stacked class-fraction bars (Dominance / Mixture / Restructuring) for the three pH-pair types in Base and Nutr$+$; pair types are ordered Acid--Alk, Acid--Acid, Alk--Alk, with Dominance \% labelled inside each bar. Reported class-frequency p-values are from two-sided Fisher exact tests on 2$\times$2 tables comparing Acid--Alk vs pooled same-pH (Acid--Acid $\cup$ Alk--Alk) events and Dominance vs non-Dominance outcomes. \textbf{(Right)} Signed PDI for acid--alk pairs: positive = acidic parental community won, negative = alkaline parental community won. Acid-win p-values are from two-sided binomial tests against the null expectation that the acidic parental community wins 50\% of acid--alk events. Dots = individual events; squares = mean $\pm$ s.e.m.; p-values above each column are from one-sample $t$-tests of signed PDI against zero.}
\end{figure}

% ============================================
\subsection*{R1-3. Circularity in Fig.~5C, PDI excluding dominant species}

\reviewercomment{In Fig.~5C, it seems like there is some risk of circularity if the dominant species are included in the calculation of the PDI. Does the positive correlation under Nutr+ conditions hold up if the dominant species are excluded from this calculation?}

\responselabel
We agree that this is an important circularity check. We used two removal strategies to avoid biasing the result toward the mixed-community winner. In the first, we removed the most abundant species in the coalesced community from all three abundance vectors; in the second, stricter test, we removed the most abundant species from each parental community separately before recalculating PDI. We followed the same PDI calculation and filtering criteria as in the original Fig.~5C analysis: events are included if they are classified as Mixture or Dominance before dominant species removal.

This analysis reveals that our core conclusion remains supported: dominant-species pairwise interactions predict the community-level winning direction in Nutr$^+$, but not in Base. In Base, $R^2$ drops from $0.11$ to $0.00 / 0.07$ and the slope from $0.49$ to $-0.02 / 0.33$ (Original / mix-winner removed / parental-community dominants removed). In Nutr$^+$, $R^2$ drops from $0.49$ to $0.24 / 0.18$ and the slope from $1.03$ to $0.64 / 0.56$ (Response Fig.~R1-3). Using independent, non-reflected events, the Spearman correlation between dominant-species pairwise score and community-level score remains significant after both removal strategies: $\rho = 0.41$ ($p = 8.2 \times 10^{-4}$) after removing the mix-winner species and $\rho = 0.42$ ($p = 1.3 \times 10^{-3}$) after removing the dominant species from both parental communities. The corresponding linear slopes are also positive and significant ($p = 1.1 \times 10^{-4}$ and $1.5 \times 10^{-4}$). Thus, the dominant species is the primary single contributor in Nutr$^+$, but the Nutr$^+$ trend does not disappear when dominant species are excluded. In contrast, Base does not retain a significant correlation after removal ($p = 0.94$ and $0.12$ for the two linear slopes; Spearman $p = 1.00$ and $0.21$).

We added to Results \S2.5 and Supplementary Fig.~31 stating: \mschange{``To test whether the Fig.~5C relationship depended solely on including the dominant species in the PDI calculation, we recalculated PDI after removing either the most abundant species in the coalesced community from all three abundance vectors or the most abundant species from each parental community before recalculating PDI, using the same Fig.~5C event filter (Mixture or Dominance before dominant-species removal). Predictive power decreased in Nutr$+$, but the association remained significant after both removal strategies (Spearman $\rho = 0.41$--$0.42$, $p \leq 1.3 \times 10^{-3}$; Supplementary Fig.~31), whereas Base did not retain a significant correlation. Thus, dominant species are major contributors to the Nutr$+$ top-down regime, but the Nutr$+$ predictability trend is not solely an artifact of including those species in the PDI calculation.''}

\begin{figure}[H]
\centering
\includegraphics[width=\textwidth]{Fig_R1_3_per_medium_scatter.pdf}
\caption*{\textbf{Response Fig.~R1-3.} \textbf{Nutr$+$ predictability persists after dominant-species removal, arguing against a PDI circularity artifact.} Per-medium PDI scatter is shown with no pooling between Base and Nutr$^+$. Rows: Base (MN, brown) and Nutr$^+$ (HN, orange). Columns: Original; Dom(mix) (zeroing argmax $c_{\text{mix}}$); Dom(parental) (zeroing argmax $c_1$ and argmax $c_2$; the stricter test). $R^2$ and slope annotations are computed on the reflected plotting data, as in Fig.~5C; any displayed slope p-values are from two-sided ordinary least-squares tests of slope $=0$. Spearman $\rho$ and $p$ are computed on the independent, non-reflected events; Spearman p-values are from two-sided rank-correlation tests of the null hypothesis of no monotonic association.}
\end{figure}

% ============================================
\subsection*{R1-4. Pool-size effects}

\reviewercomment{I wondered if the Authors could provide slightly more detail on the effect of the initial pool size. Is it similarly true that there is no significant effect of initial richness in the experimental communities? By eye, it looks like Dominance might be more likely for the 24 species pool, but it is hard to tell by squinting at Fig.~1E. It might also be interesting to see the survival ratio as a function of initial richness for the experimental communities.}

\responselabel
We do not detect a significant effect of pool size on Dominance frequency in the experiment, and we thank the reviewer for prompting this more direct check. We show this analysis separately for each nutrient condition in Response Fig.~R1-4A--C. When pooled across nutrient conditions, Dominance frequency is $61.3\%$, $58.9\%$, and $59.4\%$ for pool sizes 6, 12, and 24 species ($\chi^2 = 2.24$, $p = 0.69$; panel C), and the pool-size effect is also not significant within individual nutrient conditions (Dominance vs non-Dominance tests: Nutr$-$ $p = 0.39$, Base $p = 0.82$, Nutr$+$ $p = 0.27$). In contrast, realized parental richness clearly increases with pool size (panel A, Kruskal--Wallis $p = 1.1 \times 10^{-6}$). We also added the parental ASV retention ratio requested by the reviewer, defined as the number of parental ASV occurrences retained above threshold after coalescence divided by the summed thresholded richness of the two parental communities; this ratio is bounded by 1 and shows a modest pool-size association when pooled across media (panel B, Kruskal--Wallis $p = 0.032$). Thus, initial pool size affects richness and retention, but these effects do not translate into a significant change in Dominance frequency.

\reviewercomment{I also wondered if the Authors could look more closely at why there is no effect in the model. Perhaps by plotting the realized richness and interaction strength in parent communities as a function of initial richness?}

To address the model part of the comment, we performed a gLV pool-size ablation study at $\mu \in \{0.30, 0.60, 0.80\}$, varying the initial pool size from 4 to 48 species per parental community (100 replicate assemblies per $\mu$ and pool size). For each condition, we quantified both the Dominance frequency and the same-vs-cross pairwise-selection gap (panels E,F), and as a supporting diagnostic we reconstructed the interaction matrices from the saved simulation seeds to examine post-assembly interaction strengths. The model gives a consistent interpretation: at fixed $\mu$, Dominance frequency is comparatively insensitive to pool size (panel E), whereas the same-vs-cross pairwise-selection gap increases with pool size (panel F). Consistent with the assembly-filtering mechanism, post-assembly within-community interaction coefficients remain lower than between-community coefficients across pool sizes. This supports the interpretation that increasing pool size affects the strength of assembly filtering and correlated species fates more than it affects the categorical Dominance frequency. The experimental pairwise-selection panel added in this revision shows the same direction qualitatively: Base and Nutr$+$ show a positive same-vs-cross gap, while Nutr$-$ does not (panel D).

We added this analysis as a separate supplementary note, \mschange{``Supplementary Note 6: Pool-size dependency on community-level selection,''} and to Supplementary Fig.~32, stating: \mschange{``We tested whether initial species-pool size could account for the observed community-level selection signatures in the experiments and simulations (Supplementary Fig.~32). In the synthetic-community experiments, we did not detect a significant effect of initial pool size on Dominance frequency across 6, 12, and 24 species (61.3\%, 58.9\%, and 59.4\%, respectively; $\chi^2 = 2.24$, $p = 0.69$), and this null result held within each nutrient condition (Dominance vs non-Dominance tests: Nutr$-$ $p = 0.39$, Base $p = 0.82$, Nutr$+$ $p = 0.27$). Realized parental richness increased with initial pool size (Kruskal--Wallis $p = 1.1 \times 10^{-6}$), and parental ASV retention after coalescence varied modestly with pool size (Kruskal--Wallis $p = 0.032$), but these richness effects did not translate into a corresponding pool-size dependence of Dominance frequency.''}

\mschange{``In the gLV model, we performed a pool-size ablation at $\mu \in \{0.30, 0.60, 0.80\}$, varying the initial pool size from 4 to 48 species per parental community. Dominance frequency was primarily set by $\mu$ and changed little with pool size, whereas the same-vs-cross pairwise-selection gap increased with pool size. Reconstructed post-assembly interaction matrices showed lower within-community than between-community interaction strengths across pool sizes, supporting the interpretation that pool size affects assembly filtering and correlated species fates more than it affects Dominance-class frequency. Thus, pool-size variation does not explain the main transition in community-level selection; instead, it modulates the strength of correlated species fates within the interaction-strength regimes described in the main text.''}

\begin{figure}[H]
\centering
\includegraphics[width=\textwidth]{pool_size_analysis.pdf}
\caption*{\textbf{Response Fig.~R1-4.} \textbf{Dominance frequency is comparatively insensitive to initial pool size despite pool-size effects on richness and pairwise selection.} \textbf{A:} Realized parental richness. \textbf{B:} Parental ASV retention ratio. \textbf{C:} Experimental Dominance fraction by medium and initial pool size. \textbf{D:} Experimental pairwise species selection grouped by initial pool size. \textbf{E:} Model Dominance fraction grouped by $\mu$. \textbf{F:} Model pairwise species selection across pool sizes $4$--$48$.}
\end{figure}

% ============================================
\subsection*{R1-5. Similarity-metric robustness claim}

\reviewercomment{The Authors write (line 134) that ``This pattern of Dominance as the most frequent outcome was robust across variants of similarity metrics''. I liked the Author's approach to measuring outcomes, and I don't advocate for changing it, but this claim seems a little misleading. Two of four alternative similarity metrics gave different results, with Dominance being the \emph{least} likely outcome under Jensen--Shannon and Jaccard.}

\responselabel
We agree that our original robustness claim was too broad and could be misleading: the Dominance pattern is recovered by the alternative relative-abundance metrics Euclidean distance and Bray--Curtis dissimilarity, but not under Jaccard index or Jensen--Shannon divergence. This pattern is consistent with the fact that the metrics emphasize different aspects of community composition: relative-abundance metrics are more sensitive to quantitative dominance in relative abundance, whereas Jaccard reflects presence/absence retention and Jensen--Shannon evaluates divergence across the full abundance distribution. Thus, a coalesced community can be skewed toward one parental community in abundance space while still retaining taxa from both parental communities, or while redistributing subdominant taxa in ways that make it less skewed toward one parental community under Jaccard or Jensen--Shannon.

We revised Results \S2.1 to narrow the main-text claim to alternative relative-abundance metrics: \mschange{``This pattern of Dominance as the most frequent outcome was robust across variants of relative-abundance metrics, including Euclidean distance and Bray--Curtis dissimilarity (Extended Data Fig.~2).''} We also revised the Extended Data Fig.~2 caption to explicitly discuss the divergent metrics: \mschange{``The primary Dominance pattern is recovered under alternative relative-abundance metrics (Euclidean distance and Bray--Curtis dissimilarity), but not under Jensen--Shannon divergence or Jaccard index, for which the coalesced community is less often measured as skewed toward one parental community. This divergence reflects the fact that the metrics emphasize different aspects of community composition: relative-abundance metrics are more sensitive to quantitative abundance dominance, whereas Jaccard reflects presence/absence retention and Jensen--Shannon evaluates divergence across the full abundance distribution.''}


% ============================================
\subsection*{R1-6. Gray reflected points in figures}

\reviewercomment{It took me quite a while to realize that the gray points in Figs.~1E, 4C, 5C, and 6B are simply a reflection of the colored points. Please indicate this clearly, or perhaps just remove these.}

\responselabel
We retained the reflected points and updated the captions of Figs.~1E, 4C, 5C, and 6B to state explicitly that the gray points are reflections: \mschange{``Gray points represent the symmetric counterpart of each coalescence event (community B versus A).''}

Revised the captions of Figs.~1E, 4C, 5C, and 6B as quoted above.


% ============================================
\subsection*{R1-7. Interaction matrix after assembly}

\reviewercomment{In Fig.~2A, it might be helpful to show the matrix of interaction coefficients after assembly of the parent communities (with visible block structure).}

\responselabel
We have added the post-assembly interaction matrix as Supplementary Fig.~27, illustrating the block-diagonal structure with lower within-community than between-community competition coefficients at $\mu = 0.50$ (mean $0.389$ vs $0.500$; Mann--Whitney $U$, $p < 0.001$). This figure makes explicit the assembly-filtering structure underlying Fig.~2D.

% ============================================
\subsection*{R1-8. Fig.~2D visualization}

\reviewercomment{Fig.~2D, I was somewhat confused by the relationship between point and squares (means?). Visually, the squares do not look like they are near the means of the points. Please clarify this for readers. Additionally, the gray horizontal bars are not explained in the main text unless I missed it.}

\responselabel
The previous Fig.~2D was unclear because the individual dots and summary markers represented different display levels. In the simulation panel, the dots show a stratified subset of per-event pairwise selection correlations for visual legibility, whereas the squares and error bars show the mean $\pm$ s.e.m.\ across all simulated coalescence events. This explains why the squares need not visually coincide with the apparent center of the displayed dots. We revised the Fig.~2D caption to state this distinction explicitly and to define the gray horizontal lines as the random-selection baseline from an origin-shuffle null model: \mschange{``Individual dots show per-event pairwise selection correlations; in the simulation panel, 50 stratified events per group are displayed for legibility. Squares with error bars show mean $\pm$ s.e.m.\ across all coalescence events ($n = 1{,}200$ per group in simulations and all valid experimental events in the experimental panel). Gray horizontal lines indicate the random selection baseline from a null model in which species origin labels are shuffled (see Supplementary Information).''}


% ============================================
\subsection*{R1-9. Emphasize ``cohesion without cooperation''}

\reviewercomment{Discussion, lines 342--345: I think it is worth emphasizing how interesting and perhaps surprising it is to find ``community-level cohesion without cooperation'' as the paper by M.\ Tikhonov has it. This finding will be counterintuitive to many readers, even though the mechanism is made clear.}

\responselabel
We agree with the reviewer that this is a central and counterintuitive implication of the study. We revised the Discussion to make clear that our results are consistent with, and experimentally extend, Tikhonov's theoretical ``cohesion without cooperation'' framework. We also cite the competitive gLV literature used in our model framing (May 1972; Grilli et al.\ 2017; Hu et al.\ 2022, 2025). The revised text is placed immediately after the mechanism statement that assembly filters species into mutually weakly competing groups whose members face stronger competition from outsiders:

\mschange{``Our results are consistent with and extend Tikhonov's prediction of community-level cohesion without cooperation (Tikhonov, 2016). Whereas community-level cohesion is often associated with cooperative mechanisms such as cross-feeding or division of labor, our competition-only gLV model, building on classic random-interaction ecology (May, 1972; Grilli et al., 2017) and recent microbial gLV studies (Hu et al., 2022, 2025), shows that assembly history can generate cohesive parental units even without cooperative interactions.''}


% ============================================
\subsection*{R1-10. Means missing in Extended Data Fig.~5C}

\reviewercomment{In ED Fig.~5C, are means missing?}

\responselabel
We thank the reviewer for pointing this out. We inspected Extended Data Fig.~5C and found that the mean markers were not sufficiently visible because the panel's zoomed y-axis clipped the square mean markers. We regenerated Extended Data Fig.~5 with a wider y-axis in panel C so that the square mean markers and error bars are clearly visible. The caption already states that individual dots show per-event correlations, whereas squares with error bars show mean $\pm$ s.e.m.

We revised Extended Data Fig.~5 accordingly.


% ============================================
\subsection*{R1-11. Incorrect Extended Data reference in the SI}

\reviewercomment{On p.~9 in the SI, I believe the sentence ``Simulations were performed with parental community sizes ranging from 4 to 48 species per community (Extended Data Fig.~5)'' should refer to ED Fig.~4.}

\responselabel
We thank the reviewer for catching this error and corrected the SI reference.
